% =====================================================================
%  ABSTRACT
%
%  RULE: the abstract is written LAST. Everything it promises must exist
%  in the body. The [PENDING] marker below is where the findings go, and
%  it stays until the benchmark campaign has produced them. Project rule
%  R1: no number that was not measured.
%
%  Guidelines: 250-300 words, five keywords. `make check` measures both
%  and tells you where you stand.
% =====================================================================
\frontheading{ABSTRACT}{Abstract}

Consensus is the mechanism by which a distributed system agrees on a
single order of commands despite faults and delays. Dozens of protocols
have been proposed since 1988, and the literature has not converged on a
common basis for comparing them. Published measurements do not transfer
between papers, because they depend on the hardware, the load model and
the deployment topology. The same comparison also reverses depending on
which metric is chosen.

This thesis studies the evolution of consensus protocols and compares
selected \emph{implementations} of them in a controlled environment. It
covers the crash-fault protocols Paxos, Raft and EPaxos, and the
Byzantine protocols Tendermint and HotStuff, together with their
predecessors and descendants. Each protocol is described from its
primary source. Each move through the design space is attributed to the
specific trade-off that explains it.

The evaluation uses a single measurement instrument on a cluster of
virtual machines, with a hybrid load model. It runs open loop for
throughput and latency tails, and closed loop for the saturation search.
Predictions are registered before the measurements are taken. Every
comparative claim carries its level of comparability, so that what is
attributable to the protocol is separated from what is attributable to
implementation engineering.

\textbf{[FILL IN AFTER THE DATA: the findings, in 2--3 sentences.]}

The limits of the work are stated explicitly. What is measured is always
an implementation at a named version and configuration, not the protocol
itself. The range of cluster sizes examined experimentally is small
relative to the scales at which asymptotic behaviour appears.

\vspace{1.2em}
\noindent\textbf{SCIENTIFIC AREA:} \thScientificAreaEN.

\vspace{0.6em}
\noindent\textbf{KEYWORDS:} consensus protocols; state machine
replication; Byzantine fault tolerance; performance evaluation;
comparability.
